\begin{examplebox}

	\textbf{\textcolor{CExample}{例 1（基础）}}

	\textbf{\textcolor{CExample}{题目：}} 已知 $|\vec a|=2, |\vec b|=3$，$\vec a$ 与 $\vec b$ 的夹角 $\theta = 60^\circ$，求 $\vec a$ 在 $\vec b$ 方向上的投影数量。

	\textbf{\textcolor{CExample}{解题步骤：}}

	\begin{description}[style=nextline, leftmargin=2.8em, labelsep=0.5em, itemsep=0.45em]

		\item[\textbf{第一步（检查条件）：}] 因为 $|\vec b|=3\neq 0$，所以可以计算 $\vec a$ 在 $\vec b$ 方向上的投影数量。

		      \[
			      |\vec b|\neq 0.
		      \]

		      \par\textbf{理由：}
		      投影方向必须是非零向量。


		\item[\textbf{第二步（选用公式）：}] 已知模长和夹角，使用公式 $|\vec a|\cos\theta$。

		      \[
			      |\vec a|\cos\theta .
		      \]

		      \par\textbf{理由：}
		      此处 $\vec a,\vec b$ 均为非零向量，夹角公式可以直接使用。


		\item[\textbf{第三步（代入数据）：}] 代入 $|\vec a|=2$，$\cos60^\circ=\dfrac12$。

		      \[
			      2 \cdot \frac12 = 1 .
		      \]

		      \par\textbf{理由：}
		      替换已知数值并计算。

	\end{description}

	\textbf{\textcolor{CExample}{关键结论：}} $\vec a$ 在 $\vec b$ 方向上的投影数量为 $\mathbf{1}$。

	\medskip

	\textbf{\textcolor{CExample}{例 2（坐标计算）}}

	\textbf{\textcolor{CExample}{题目：}} 已知向量 $\vec a=(1,2)$，$\vec b=(3,4)$，求 $\vec a$ 在 $\vec b$ 方向上的投影数量。

	\textbf{\textcolor{CExample}{解题步骤：}}

	\begin{description}[style=nextline, leftmargin=2.8em, labelsep=0.5em, itemsep=0.45em]

		\item[\textbf{第一步（计算点积）：}] 坐标点积公式为 $\vec a\cdot\vec b=x_1x_2+y_1y_2$。

		      \[
			      \vec a\cdot\vec b
			      =
			      1\cdot 3+2\cdot 4
			      =
			      11.
		      \]

		      \par\textbf{理由：}
		      将坐标代入点积定义。


		\item[\textbf{第二步（计算模长）：}] 计算 $\vec b$ 的模长。

		      \[
			      |\vec b|
			      =
			      \sqrt{3^2+4^2}
			      =
			      5.
		      \]

		      \par\textbf{理由：}
		      投影公式的分母是投影方向向量 $\vec b$ 的模长。


		\item[\textbf{第三步（套用投影公式）：}] 投影数量为

		      \[
			      \frac{\vec a\cdot\vec b}{|\vec b|}
			      =
			      \frac{11}{5}.
		      \]

		      \par\textbf{理由：}
		      投影数量是标量，不是向量。

	\end{description}

	\textbf{\textcolor{CExample}{关键结论：}} $\vec a$ 在 $\vec b$ 方向上的投影数量为 $\mathbf{\dfrac{11}{5}}$。

	\medskip

	\textbf{\textcolor{CExample}{例 3（综合应用）}}

	\textbf{\textcolor{CExample}{题目：}} 已知 $\vec a$ 在 $\vec b$ 方向上的投影数量为 $3$，且 $|\vec b|=2$，求 $\vec a\cdot\vec b$。

	\textbf{\textcolor{CExample}{解题步骤：}}

	\begin{description}[style=nextline, leftmargin=2.8em, labelsep=0.5em, itemsep=0.45em]

		\item[\textbf{第一步（回忆公式）：}] 投影数量公式为

		      \[
			      \frac{\vec a\cdot\vec b}{|\vec b|}.
		      \]

		      \par\textbf{理由：}
		      由结论知，投影数量等于点积除以投影方向向量的模长。


		\item[\textbf{第二步（代入已知）：}] 将投影数量 $3$ 和 $|\vec b|=2$ 代入公式。

		      \[
			      3
			      =
			      \frac{\vec a\cdot\vec b}{2}.
		      \]

		      \par\textbf{理由：}
		      题目已给出投影数量和 $\vec b$ 的模长。


		\item[\textbf{第三步（求解点积）：}] 两边乘以 $2$，得

		      \[
			      \vec a\cdot\vec b
			      =
			      6.
		      \]

		      \par\textbf{理由：}
		      等式两边同乘非零数 $2$。

	\end{description}

	\textbf{\textcolor{CExample}{关键结论：}} $\vec a\cdot\vec b=\mathbf{6}$。

	\medskip

	\textbf{\textcolor{CExample}{例 4（负投影提醒）}}

	\textbf{\textcolor{CExample}{题目：}} 已知 $|\vec a|=4$，$\vec a$ 与 $\vec b$ 的夹角为 $120^\circ$，求 $\vec a$ 在 $\vec b$ 方向上的投影数量。

	\textbf{\textcolor{CExample}{解题步骤：}}

	\begin{description}[style=nextline, leftmargin=2.8em, labelsep=0.5em, itemsep=0.45em]

		\item[\textbf{第一步（选用夹角公式）：}] 投影数量为

		      \[
			      |\vec a|\cos\theta.
		      \]

		      \par\textbf{理由：}
		      已知 $|\vec a|$ 和夹角，可直接使用夹角形式。


		\item[\textbf{第二步（代入计算）：}] 因为 $\cos120^\circ=-\dfrac12$，所以

		      \[
			      |\vec a|\cos120^\circ
			      =
			      4\cdot\left(-\frac12\right)
			      =
			      -2.
		      \]

		      \par\textbf{理由：}
		      夹角为钝角时，余弦值为负，投影数量也为负。

	\end{description}

	\textbf{\textcolor{CExample}{关键结论：}} $\vec a$ 在 $\vec b$ 方向上的投影数量为 $\mathbf{-2}$，表示投影方向与 $\vec b$ 方向相反。

\end{examplebox}